\clearpage\SourceLeaf{441}\ArithmeticHead{434}{L}{ANALYSIS. BRIEFWECHSEL.}
\NumberedItem{5}
\HistoricTitle{\textsc{Gauss} an \textsc{Schumacher}.\quad Göttingen, 1. September 1850.}
\HistoricCitation{[Briefwechsel zwischen \textsc{Gauss} und \textsc{Schumacher} VI, Altona 1865, S. 107.]}
\noindent \(\,{.}\;{.}\;{.}\;{.}\) Indem ich eben im Begriff war, diesen Brief zu schliessen, erhielt ich die von Ihnen übersandte \textsc{Prehn}sche Abhandlung[\textsuperscript{*})], von der ich doch erst flüchtig Einsicht nehmen wollte. Darf ich mich offen gegen Sie aussprechen, so glaube ich nicht, dass eine Reise des H[errn] Verfassers zu seiner Zufriedenheit ausfallen würde. Meine Anforderungen an mathematische Argumentation in Beziehung auf Strenge und Klarheit liegen in der That von denen des Verfassers soweit entfernt, dass ich an einer gegenseitigen Verständigung zweifle; wir stehen auf ganz verschiedenem Boden, wie Conservativer und Radicaler. Solche Gleichnisse hinken allerdings: in \textls[120]{diesem} Falle wäre ich der Radicale, der kein historisches Recht anerkennt, sondern nichts ohne strengen Beweis des Rechtstitel[s] gelten lässt. Der Unterschied ist aber, dass im Leben Conservativer und Radicaler ihre Ansprüche à tout prix realisiren, in der Wissenschaft aber die strengsten Forderungen nur zu unserer eigenen Befriedigung gemacht werden, und man den andern, der sich mit laxen oder unklarn Beweisen begnügen mag, gern gewähren lässt. Der eigentliche Kern der Sache ist wie mir deucht folgender.

\indent Es ist der Character der Mathematik der neueren Zeit (im Gegensatz gegen das Alterthum), dass durch unsere Zeichensprache und Namengebungen wir einen Hebel besitzen, wodurch die verwickeltsten Argumentationen auf einen gewissen Mechanismus reducirt werden. An Reichthum hat dadurch die Wissenschaft unendlich gewonnen, an Schönheit und Solidität aber, wie das Geschäft gewöhnlich betrieben wird, eben so sehr verloren. Wie oft wird jener Hebel eben nur mechanisch angewandt, obgleich die Befugniss dazu in den meisten Fällen gewisse stillschweigende Voraussetzungen implicirt. Ich fordere, man soll bei allem Gebrauch des Calculs, bei allen Begriffsverwen-
\par\medskip\noindent\rule{.24\linewidth}{.4pt}\par\smallskip
{\small\noindent\hspace*{\parindent}[\textsuperscript{*}) \textsc{Jeppe Prehn}, \textit{Über die Bedeutung der divergenten unendlichen Reihen, die Bestimmung ihrer Werthe, und über die Zulässlichkeit ihrer Anwendung bei analytischen Rechnungen}, \textsc{Crelles} Journal für Mathematik \textbf{41} (1850), S. 1, vergl. \textit{Berichtigungen}, ebenda S. 364, die letzteren sind nach dem am 28. November 1850 erfolgten Tode des Verfassers erschienen.]\par}
\clearpage\SourceLeaf{442}\ArithmeticHead{435}{R}{GAUSS AN SCHUMACHER.}
\noindent dungen sich immer der ursprünglichen Bedingungen bewusst bleiben, und \textls[120]{alle} Producte des Mechanismus niemals über die klare Befugniss hinaus als Eigenthum betrachten. Der gewöhnliche Gang ist aber der, dass man für die Analysis einen Character der Allgemeinheit in Anspruch nimmt, und dem \textls[120]{Andern}, der so herausgebrachte Resultate noch nicht für bewiesen anerkennt, zumuthet, er solle das Gegentheil nachweisen. Diese Zumuthung darf man aber nur an den stellen, der seinerseits \textls[120]{behauptet} ein Resultat sei falsch, nicht aber dem, der ein Resultat nicht für bewiesen anerkennt, welches auf einem Mechanismus beruhet, dessen ursprüngliche, wesentliche Bedingungen in dem vorliegenden Fall gar nicht zutreffen. So ist es sehr oft mit divergirenden Reihen. Reihen haben eine klare Bedeutung, wenn sie convergiren; diese Klarheit der Bedeutung fällt weg mit dieser Bedingung, und es ändert im Wesentlichen Nichts, ob man sich des Worts Summe oder Werth bedient. Der Raum eines Briefes ist aber viel zu klein, um alles weiter auszuführen. — Nehmen Sie meinetwegen statt obigen Gleichnisses einer Maschine das von Papiergeld. Es kann dies zu grossen Arbeiten vortheilhaftest benutzt werden, aber solide ist der Gebrauch nur, wenn ich gewiss bin, es jeden Augenblick in klingende Münze umsetzen zu können.

\indent Es scheint mir übrigens aus einigen Indicien hervorzugehen, dass der Verfasser mit \textls[120]{meinen} Arbeiten nicht bekannt ist. Er meint die Gültigkeit des Gebrauchs der divergenten Reihen sei allgemein unbedenklich anerkannt noch in den ersten Decennien des gegenwärtigen Jahrhunderts. Ich habe sie nie anerkannt; zwar niemals ex professo dagegen geschrieben, aber überall, wo eine Veranlassung war, die Zulässigkeit der Reihen nur unter der Bedingung der Convergenz als sich von selbst verstehend entschieden ausgesprochen. In diesem Augenblick würde ich nur [hinweisen] z. B. auf meine Schrift von 1799, p. 12[\textsuperscript{*})]. Meine Schrift von 1812 über die Reihe \(1+\frac{\alpha\beta}{1\,{.}\,\gamma}x\) \textit{etc.}[\textsuperscript{**})]. Die Anzeige meiner Schrift über die Anziehung der ellipt[ischen] Sphäroide in den Göttinger gelehrten Anzeigen 1813, p. 547, auch abgedruckt Monatliche Correspondenz 27. Bd. p. 424[\textsuperscript{***})]. Auch die Art wie H. \textsc{Pr[ehn]} über die
\par\medskip\noindent\rule{.24\linewidth}{.4pt}\par\smallskip
{\small
\noindent\hspace*{\parindent}[\textsuperscript{*}) Werke III, S. 9—10, Absatz 3.]
\par\noindent\hspace*{\parindent}[\textsuperscript{**}) Werke III, S. 123 ff.]
\par\noindent\hspace*{\parindent}[\textsuperscript{***}) Werke V, S. 281 oben.]
\par}
\par\smallskip\noindent\hfill{\small 55\textsuperscript{*}}
\clearpage\SourceLeaf{443}\ArithmeticHead{436}{L}{ANALYSIS. BRIEFWECHSEL.}
\noindent imaginären Grössen spricht, zeigt, dass er sich noch ganz auf dem Standpunkte befindet, auf dem man sich vor 1831 befand, und die gänzlich veränderte Gestalt dieser Lehre, die ich ihr gegeben habe, nur auf ein Paar Seiten, aber den Kern der Sache erschöpfend, gar nicht kennt. \(\,{.}\quad{.}\quad{.}\quad{.}\quad{.}\)
